
\newcommand{\betamin}{\beta_{\textsc{min}}}

\subsubsection{Frequentist consistency metric}
\label{cfc/app:frequentist-metric}
\todo{maybe change P to x throughout here too}
In a deterministic world, we cannot let any inconsistency pass; every time we prove any rule of probability does not hold exactly, we must discard the forecaster as flawed.
This is too strict for the consistency check framework to be useful.
Instead, we propose a violation metric and the corresponding inconsistency threshold based on statistical hypothesis testing.\todo{make equations break properly}

Assume that each event $P$ has a true probability value $\TrueProb(P)$, say under some world model that accounts for aleatoric uncertainty.
\begin{definition}[Frequentist consistency]
\label{def:frequentist-consistency}
A frequentist-consistent forecaster $\Forecaster$ samples a Gaussian estimate $\TrueProb(P) + \eps$ of each event $P$, with variance $\sigma^2 \TrueProb(P) (1 - \TrueProb(P))$ for a hyperparameter $\sigma^2$:
\begin{align}
    \label{eq:frequentist_consistency}
    \Forecaster(P) - \TrueProb(P) \sim \Normal\left(0, \sigma^2 \TrueProb(P) (1 - \TrueProb(P))\right) \quad \text{independently for all events } P.
\end{align}
\end{definition}
This is principled from the frequentist perspective.
Consider a forecaster that just samples the (relevant subset of) the world $n$ times using the best available world simulator, and estimates the probability of each event $P$ as the proportion of times that $P$ occurs in the $n$ samples.
If we estimate the probability as the average chance of an event $P$ with true probability $p$ occurring out of $n$ times, then this estimate has a scaled binomial distribution with mean $p$ and variance $p(1-p) / n$.
To reach \Cref{eq:frequentist_consistency}, replace the averaged binomial with the Gaussian of the same variance, and denote $\sigma^2 := 1/n$.


This simple model enables us to derive hypothesis tests for each of the consistency checks described in \Cref{tab:consistency_checks}.
The null hypothesis is always that the forecaster is frequentist-consistent.
Note that $\sigma^2$ is not our estimate of the variance of any forecaster; it is just a hyperparameter that controls how strict our null hypothesis is.
We leave estimating the variance of a particular forecaster and testing frequentist consistency based on that alone to future work.

\paragraph{Notation}
The expression $a \Normal(0, c^2)$ denotes a Gaussian random variable with mean 0 and variance $a^2 c^2$.
The expression $a \Normal(0, c^2) + b \Normal(0, c^2)$ denotes a Gaussian random variable with mean 0 and variance $a^2 c^2 + b^2 c^2$.
All sums range over the cyclic permutations of the variables under the sum.
All $\Normal(0, c^2)$ terms appearing with the same power of $\sigma$ are independent.
Two $\Normal(0, c^2)$ terms appearing with a different power of $\sigma$ may be correlated; this is not important for our purposes, since we discard high-order powers of $\sigma$.

\paragraph{Bootstrapping the true probability}
The final expressions for hypothesis test statistics might involve the true probability $\TrueProb(P)$.
It is not available, so we just plug in $\Forecaster(P)$ for $\TrueProb(P)$ in the end.
If we had a prior on $\TrueProb(P)$, we could combine it with $\Forecaster(P)$ to get a more robust estimate.


\paragraph{\NegChecker}
We take the violation metric and the corresponding threshold as to produce a hypothesis test against this:
\begin{align*}
   \Forecaster(P) + \Forecaster(\neg P) - 1
   = \TrueProb(P) + \eps_1 + \TrueProb(\neg P) + \eps_2 - 1
   = \eps_1 + \eps_2
   \\
   \sim \Normal\left(0, \sigma^2 (\TrueProb(P)(1 - \TrueProb(P)) + \TrueProb(\neg P)(1 - \TrueProb(\neg P))) \right)
\end{align*}


We estimate the unknown $\TrueProb$ values with the corresponding $\Forecaster$ estimates.  Note that, although $\TrueProb(P) = 1- \TrueProb(\neg P)$, it is of course not necessarily the case that $\Forecaster(P) = 1- \Forecaster(\neg P)$.

The error distribution is $\sigma \Normal\left( \Forecaster(P) (1 - \Forecaster(P)) + \Forecaster(\neg P) (1 - \Forecaster(\neg P)) \right)$, and the two-sided test is
\begin{align*}
   \abs{\Forecaster(P) + \Forecaster(\neg P) - 1} < \gamma \sigma \sqrt{(1 - \Forecaster(P)) \Forecaster(P) + (1 - \Forecaster(\neg P)) \Forecaster(\neg P)}
\end{align*}
for some scale factor $\gamma$ (number of standard deviations) that scales the power of the test. For example, $\gamma = 2.58$ gives a 99\%-confidence interval.

We now want to compute some \emph{consistency violation metric} that makes inconsistency comparable across different checks. The natural idea is to aggregate all terms dependent on $\Forecaster$ to one side; and make the hypothesis test be just some threshold on the computed violation metric.

It is possible that the denominator of the resulting expression is $0$ when the forecaster is certain and $\Forecaster$ is $0$ or $1$; to avoid division with zero, we add a small regularization term $\betamin = 10^{-3}$.
See the last paragraph of this section for a discussion of hyperparameters.

Our consistency violation metric is then:
\begin{align*}
    v_{\NegChecker}
    = \frac{\abs{\Forecaster(P) + \Forecaster(\neg P) - 1}}{\sqrt{(1 - \Forecaster(P)) \Forecaster(P) + (1 - \Forecaster(\neg P)) \Forecaster(\neg P) + \betamin}}.
\end{align*}

The hyperparameter $\sigma^2$ determines how strict we are with rejecting inconsistencies which could be attributed to ``noisy'' predictions. Note that the violation metric itself does not depend on $\sigma^2$.

A violation (inconsistency), therefore, occurs when:
\begin{align*}
    v_{\NegChecker} > \gamma \sigma.
\end{align*}

\paragraph{\CondCondChecker}
This is a more complex consistency check; we derive the hypothesis test and violation metric in detail below. For the other checks, we just report the short derivation.

\begin{align*}
(a, b, c, d) &= \left(\TrueProb(P), \TrueProb(Q \mid P), \TrueProb(R \mid P \wedge Q), \TrueProb(P \wedge Q \wedge R)\right) \\
(a', b', c', d') &= \left(\Forecaster(P), \Forecaster(Q \mid P), \Forecaster(R \mid P \wedge Q), \Forecaster(P \wedge Q \wedge R)\right)
\end{align*}

We can write:

\begin{align*}
    \Forecaster(P) &= \Normal\left(0, \sigma^2 a (1 - a)\right) + a, \\
    \Forecaster(Q\mid P) &= \Normal\left(0, \sigma^2 b (1 - b)\right) + b, \\
    \Forecaster(R \mid P \wedge Q) &= \Normal\left(0, \sigma^2 c (1 - c)\right) + c, \\
    \Forecaster(P \wedge Q \wedge R) &= \Normal\left(0, \sigma^2 d (1 - d)\right) + d
\end{align*}

We now compute the difference of the two expressions that should be equal.
All sums and products are cyclic over $a$, $b$, $c$.
%, we can compute where $i,j,k$ cyclically take on the values of $a,b,c$
%\begin{align*}
%    \Forecaster(P) \Forecaster(Q \mid P) \Forecaster(R \mid P \wedge Q) - \Forecaster(P \wedge Q \wedge R) &= abc - d \\
%    &+ \sigma \left( \sum jk \Normal(0, i(1-i)) - \Normal(0, d(1-d))\right) \\
%    &+ \sigma^2 \sum i \Normal(0, j(1-j)) \Normal(0, k(1-k)) \\
%    &+ \sigma^3 \prod \Normal(0, i(1-i))
%\end{align*}
\begin{align*}
    \Forecaster(P) \Forecaster(Q \mid P) \Forecaster(R \mid P \wedge Q) - \Forecaster(P \wedge Q \wedge R)
    &= abc - d \\
    &+ \sigma \left( \sum_{a} bc \Normal(0, a(1-a)) - \Normal(0, d(1-d))\right) \\
    &+ \sigma^2 \sum_{a} \Normal(0, b(1-b)) \Normal(0, c(1-c)) \\
    &+ \sigma^3 \prod_{a} \Normal(0, a(1-a)).
\end{align*}

In the above, all Gaussians with the same variance are identical, and all other combinations are independent.
As $abc - d = 0$ by the law of total probability, the leading error term is next to $\sigma$.
This is a Gaussian with mean 0 and standard deviation:

\begin{align*}
    \sigma \sqrt{\sum_{a} b^2c^2 a(1-a) + d(1-d)}
    =
    \sigma \sqrt{abc \sum_{a} bc (1-a) + d(1-d)}
\end{align*}
\todo{this expression could be simplified, it's already simplified in code}



We now discard the terms of $\sigma^2$, $\sigma^3$, and in general any higher order power of $\sigma$. This is principled because the coefficients can always be (in some confidence interval) upper bounded by a constant independent of $\sigma$.
Hence, if $\sigma$ is small enough, the resulting test will be very close to the true hypothesis test.



We do not have the true probabilities $a$, $b$, $c$, $d$, so we just plug in $(a', b', c', d') = (\Forecaster(P), \Forecaster(Q \mid P), \Forecaster(R \mid P \wedge Q), \Forecaster(P \wedge Q \wedge R))$. \footnote{Depending on how we use the relation $abc=d$, we can end up with different expressions in the end. We choose the one that, after plugging in, (i) yields an expression for variance that is always nonnegative, and (ii) is not a polynomial multiple of any single value of $\Forecaster$.}
 Thus the hypothesis test is (where the sum is cyclic over $a'$, $b'$, $c'$):

\begin{align*}
    \abs{a'b'c' - d'}
    > \gamma \sigma \sqrt{ a' b' c' \sum_{a'} b'c' (1-a') + d' (1-d')}
\end{align*}

Our violation metric is then:
\begin{align*}
    v_{\CondCondChecker}
    = \frac{
        \abs{a'b'c' - d'}
        }{
        \sqrt{a'b'c' \sum_{a'} b'c' (1-a') + d' (1-d')+ \betamin}
    }.
\end{align*}
where again $(a', b', c', d') = (\Forecaster(P), \Forecaster(Q \mid P), \Forecaster(R \mid P \wedge Q), \Forecaster(P \wedge Q \wedge R))$ are the forecasts.


\paragraph{\CondChecker}

Similarly as for \CondCondChecker: we denote
$(a, b, c) = \left(\TrueProb(P),  \TrueProb(P \mid Q), \TrueProb(P \wedge Q) \right)$ and the associated $(a', b', c')$ for the forecasts.
Then we can compute
\begin{align*}
    &\Forecaster(P) \Forecaster(Q \mid P) - \Forecaster(P \wedge Q) \\
    &= ab - c + \sigma \left(b \Normal(0, a(1-a)) + a \Normal(0, b(1-b)) - \Normal(0, c(1-c)) \right) \\
    &\quad + \sigma^2 \Normal(0, a(1-a)) \Normal(0, b(1-b)).
\end{align*}

% \begin{align*}
%     &\Forecaster(P) \Forecaster(Q \mid P) - \Forecaster(P \wedge Q)
%     \\ &= ab - c
%     + \sigma \left(b \Normal(0, a(1-a) + a \Normal(0, b(1-b)) - \Normal(0, c(1-c)) \right)
%     + \sigma^2 \Normal(0, a(1-a)) \Normal(0, b(1-b)).
% \end{align*}



The term next to $\sigma$ is a Gaussian with mean 0 and standard deviation:

\begin{align*}
    \sigma \sqrt{a^2 b(1-b) + b^2 a(1-a) + c(1-c)}
    = \sigma \sqrt{ab \left( a(1-b) + b(1-a) \right) + c(1-c)}.
\end{align*}

Again, we have to plug in $(a', b', c') = (\Forecaster(P), \Forecaster(Q \mid P), \Forecaster(P \wedge Q))$ instead of $(a, b, c)$.

Our violation metric is then:
\begin{align*}
    v_{\CondChecker}
    = \frac{\abs{a'b' - c'}}{
        \sqrt{a'b' \left( a'(1-b') + b'(1-a') \right) + c'(1-c')+ \betamin}
    }
\end{align*}
And the test is again, for a suitable $\gamma$ corresponding to the desired power of the test:
\begin{align*}
    v_{\CondChecker} > \gamma \sigma.
\end{align*}


\paragraph{\ParaphraseChecker}

Here we can simply check whether $P$ and $Q$ are the same.

\begin{align*}
   &\Forecaster(P) - \Forecaster(Q)
   = \TrueProb(P) + \eps_1 - \TrueProb(Q) - \eps_2 \\
   &= \eps_1 - \eps_2 \sim \Normal\left(0, \sigma^2 ((\TrueProb(P)(1 - \TrueProb(P)) + (\TrueProb(Q)(1 - \TrueProb(Q)) \right)
\end{align*}

This yields the following violation metric:

\begin{align*}
    v_{\ParaphraseChecker}
    = \frac{\abs{\Forecaster(P) - \Forecaster(Q)}}{
        \sqrt{(\Forecaster(P)(1 - \Forecaster(P)) + (\Forecaster(Q)(1 - \Forecaster(Q))+ \betamin}
    }
\end{align*}

\paragraph{\AndOrChecker}

\begin{align*}
   &\Forecaster(P) + \Forecaster(Q) - \Forecaster(P \lor Q) - \Forecaster(P \land Q) \\
   & =
   \TrueProb(P) + \TrueProb(Q) - \TrueProb(P \lor Q) - \TrueProb(P \land Q) + \eps_1 + \eps_2 - \eps_3 - \eps_4 \\
   &= \eps_1 + \eps_2 - \eps_3 - \eps_4 \\
   &\sim \Normal\left(0, \sigma^2 \left( \TrueProb(P)(1 - \TrueProb(P)) + \TrueProb(Q)(1 - \TrueProb(Q)) \right. \right. \\
   &\left. \left. + \TrueProb(P \lor Q)(1 - \TrueProb(P \lor Q)) + \TrueProb(P \land Q)(1 - \TrueProb(P \land Q)) \right) \right).
\end{align*}


% \begin{align*}
%    &\Forecaster(P) + \Forecaster(Q) - \Forecaster(P \lor Q) - \Forecaster(P \land Q) \\
%    & =
%    \TrueProb(P) + \TrueProb(Q) - \TrueProb(P \lor Q) - \TrueProb(P \land Q) + \eps_1 + \eps_2 - \eps_3 - \eps_4 \\
%    &= \eps_1 + \eps_2 - \eps_3 - \eps_4 \\
%    &\sim \Normal\left(0, \sigma^2 ((\TrueProb(P)(1 - \TrueProb(P)) + (\TrueProb(Q)(1 - \TrueProb(Q)) \\
%    &+ (\TrueProb(P \lor Q)(1 - \TrueProb(P \lor Q)) + (\TrueProb(P \land Q)(1 - \TrueProb(P \land Q))) \right).
% \end{align*}

We again plug in $\Forecaster$ instead of $\TrueProb$ to compute the error term allowed: $\gamma \sigma \sqrt{M}$ where

\begin{align*}
    &M = \Forecaster(P)(1 - \Forecaster(P)) + \Forecaster(Q)(1 - \Forecaster(Q) + \Forecaster(P \lor Q)(1 - \Forecaster(P \lor Q)) + \\
    &\Forecaster(P \land Q)(1 - \Forecaster(P \land Q))
\end{align*}

and violation metric:

\begin{align*}
    v_{\AndOrChecker}
    = \frac{\abs{\Forecaster(P) + \Forecaster(Q) - \Forecaster(P \lor Q) - \Forecaster(P \land Q)}}{
        \sqrt{\begin{array}{l}
            \Forecaster(P)(1 - \Forecaster(P)) + \Forecaster(Q)(1 - \Forecaster(Q)) + \\
            \Forecaster(P \lor Q)(1 - \Forecaster(P \lor Q)) + \Forecaster(P \land Q)(1 - \Forecaster(P \land Q)) + \betamin
        \end{array}}
    }.
\end{align*}

\paragraph{\ButChecker}

\begin{align*}
   &\Forecaster(P \lor Q) - \Forecaster(P) - \Forecaster(\lnot P \land Q) =
   \TrueProb(P \lor Q) - \TrueProb(P) - \TrueProb(\lnot P \land Q) + \eps_1 - \eps_2 - \eps_3 =\\
   &\eps_1 - \eps_2 - \eps_3 \sim\\
   &\Normal\left(0, \sigma^2 ((\TrueProb(P \lor Q)(1 - \TrueProb(P \lor Q)) + (\TrueProb(P)(1 - \TrueProb(P)) + (\TrueProb(\lnot P \land Q)(1 - \TrueProb(\lnot P \land Q)) \right)
\end{align*}

with error term:

\begin{align*}
    \gamma \sigma \sqrt{\Forecaster(P \lor Q)(1 - \Forecaster(P \lor Q) + \Forecaster(P)(1 - \Forecaster(P) + \Forecaster(\lnot P \land Q)(1 - \Forecaster(\lnot P \land Q)}
\end{align*}

and violation metric:

\begin{align*}
    v_{\ButChecker}
    = \frac{\abs{\Forecaster(P \lor Q) - \Forecaster(P) - \Forecaster(\lnot P \land Q)}}{
        \sqrt{\Forecaster(P \lor Q)(1 - \Forecaster(P \lor Q)) + \Forecaster(P)(1 - \Forecaster(P)) + \Forecaster(\lnot P \land Q)(1 - \Forecaster(\lnot P \land Q)+ \betamin}
    }
\end{align*}


\paragraph{\ConsequenceChecker}

In the case of inequalities involving $\le$, there are two ways in which the consistency check can be passed.
If $\Forecaster(P) \le \Forecaster(Q)$, the consistency check is automatically passed.
Otherwise,
we check for pseudo-equality using the same violation metric as in \ParaphraseChecker.

\begin{align*}
    v_{\ConsequenceChecker}
    = [\Forecaster(P) > \Forecaster(Q)] \frac{\abs{\Forecaster(P) - \Forecaster(Q)}}{
        \sqrt{\Forecaster(P)(1 - \Forecaster(P)) + \Forecaster(Q)(1 - \Forecaster(Q))+ \betamin}
    }
\end{align*}
where $[\Forecaster(P) > \Forecaster(Q)]$ is the Iverson Bracket (1 if true, 0 otherwise).

%\textbf{Note: Adam  Not sure if we should include changing $\gamma = 2$}

%The only difference is that the hypothesis test is one-sided; therefore, $\gamma$ needs to be smaller to get the same confidence interval.
%We choose $\gamma = 2$ by default for one-sided tests, to get a level of around 97\% similar as for $\gamma = 3$ in \ParaphraseChecker.

\paragraph{\AndChecker}


Similarly to \ConsequenceChecker, if the chain of strict inequalities
\begin{align*}
    \max(\Forecaster(P) + \Forecaster(Q) - 1, 0) < \Forecaster(P \land Q) < \min(\Forecaster(P), \Forecaster(Q))
\end{align*}
holds, then the check automatically passes.  We set $v_{\AndChecker\_\textsc{lhs}} = 0$ and $v_{\AndChecker\_\textsc{rhs}} = 0$ if it passes the first and second strict inequality respectively.

If not, then we test for pseudo-equality for the violating pair:

$\text{LHS}: \max(\Forecaster(P)+\Forecaster(Q)-1, 0) = \Forecaster(P\land Q)$

$ \text{RHS}: \Forecaster(P\land Q)  = \min(\Forecaster(P),\Forecaster(Q))$

Equality check if it fails the first inequality:

\begin{align*}
    \eps_{\textsc{lhs}} = \begin{dcases*}
        \gamma\sigma
        \sqrt{\Forecaster(P)(1-\Forecaster(P))+\Forecaster(Q)(1-\Forecaster(Q)) + \Forecaster(P\land Q)(1-\Forecaster(P\land Q))} \\
        \quad \text{if } \Forecaster(P)+\Forecaster(Q)-1 > 0, \\[6pt]
        \text{N/A} \\
        \quad \text{otherwise pass as } \Forecaster(P\land Q) \ge 0.
    \end{dcases*}
\end{align*}





\begin{align*}
    v_{\AndChecker\_\textsc{lhs}}
    &= [\Forecaster(P) + \Forecaster(Q) - 1 > \Forecaster(P\land Q)] \cdot \\ &\frac{
            \Forecaster(P) + \Forecaster(Q) - 1 - \Forecaster(P \land Q)
    }{
        \sqrt{\Forecaster(P)(1-\Forecaster(P))+\Forecaster(Q)(1-\Forecaster(Q)) + \Forecaster(P\land Q)(1-\Forecaster(P\land Q))+ \betamin}
    }
\end{align*}

Equality check if it fails the second inequality:

Define $\Forecaster(R) =  \min(\Forecaster(P),\Forecaster(Q))$.

\begin{align*}
    \eps_{\textsc{rhs}}
    &=  \gamma\sigma\sqrt{
        \Forecaster(P\land Q)(1-\Forecaster(P\land Q))+\Forecaster(R)(1+\Forecaster(R))
    }
    \\
    v_{\AndChecker\_{\textsc{rhs}}}
    &= [\Forecaster(R) < \Forecaster(P\land Q)] \frac{
            \Forecaster(P \land Q) - \Forecaster(R)
    }{
        \sqrt{\Forecaster(P\land Q)(1-\Forecaster(P\land Q)) + \Forecaster(R)(1-\Forecaster(R))+ \betamin}
    }
\end{align*}


Consistency is violated if either inequality is violated, \emph{and} the respective hypothesis test for pseudo-equality fails. We use $v_{\AndChecker\_\textsc{lhs}}$ for the first and $v_{\AndChecker\_\textsc{rhs}}$ for the second inequality.
We define $v_{\AndChecker} = \max\{v_{\AndChecker\_\textsc{lhs}}, v_{\AndChecker\_\textsc{rhs}} \}$.



\paragraph{\OrChecker}

We proceed similarly as for \AndChecker{}.

If the strict inequality $\max(\Forecaster(P),\Forecaster(Q))<\Forecaster(P\lor Q)<\min(1,\Forecaster(P)+\Forecaster(Q))$ holds, then it automatically passes.  We set $v_{\OrChecker\_\textsc{lhs}} = 0$ and $v_{\OrChecker\_\textsc{rhs}} = 0$ if it passes the first and second strict inequality respectively.

If not, we test for pseudo-equality:

$\text{LHS}: \max(\Forecaster(P),\Forecaster(Q)) = \Forecaster(P\lor Q)$

$\text{RHS}: \Forecaster(P\lor Q) = \min(1,\Forecaster(P)+\Forecaster(Q))$.


Equality check LHS:
Define  $\Forecaster(S) =  \max(\Forecaster(P),\Forecaster(Q))$.
\begin{align*}
    &\eps_{\textsc{lhs}} =
        \gamma\sigma\sqrt{\Forecaster(S)(1-\Forecaster(S)) + \Forecaster(P\lor Q)(1-\Forecaster(P\lor Q))}
\end{align*}

\begin{align*}
    v_{\OrChecker\_{\textsc{lhs}}}
    = [\Forecaster(S) > \Forecaster(P\lor Q)] \frac{
      \Forecaster(S) - \Forecaster(P \lor Q)
    }{
      \sqrt{\Forecaster(S)(1-\Forecaster(S)) + \Forecaster(P\lor Q)(1-\Forecaster(P\lor Q)) + \betamin}
    }
\end{align*}

Equality check RHS:

\begin{align*}
    \eps_{\textsc{rhs}} = \begin{dcases*}
        \gamma\sigma\sqrt{\Forecaster(P\lor Q)(1-\Forecaster(P\lor Q)) + \Forecaster(P)(1-\Forecaster(P)) + \Forecaster(Q)(1-\Forecaster(Q))} \\
        \quad \text{if } \Forecaster(P)+\Forecaster(Q) < 1, \\[6pt]
        \text{N/A} \\
        \quad \text{otherwise pass as } \Forecaster(P\lor Q) \le 1.
    \end{dcases*} \\
\end{align*}
\begin{align*}
    v_{\OrChecker\_{\textsc{rhs}}}
    &= [\Forecaster(P) + \Forecaster(Q) < \Forecaster(P\lor Q)] \cdot \\ & \quad{} \frac{
            \Forecaster(P\lor Q) - \Forecaster(P) - \Forecaster(Q)
    }{
        \sqrt{\Forecaster(P\lor Q)(1-\Forecaster(P\lor Q)) + \Forecaster(P)(1-\Forecaster(P)) + \Forecaster(Q)(1-\Forecaster(Q))+ \betamin}
    }
\end{align*}

Consistency is violated if either inequality is violated, \emph{and} the subsequent hypothesis test for pseudo-equality fails. We use $v_{\OrChecker\_\textsc{lhs}}$ for the first and  $v_{\OrChecker\_\textsc{rhs}}$ for the second inequality.
Analogously to \AndChecker{}, define $v_{\OrChecker} = \max\{v_{\OrChecker\_\textsc{lhs}}, v_{\OrChecker\_\textsc{rhs}} \}$.

\paragraph{\ExpectedEvidenceChecker}
\label{cfc/app:frequentist-ee}

Write $(a,b,c,d)=(\TrueProb(P),\TrueProb(P\mid Q),\TrueProb(P\mid\lnot Q), \TrueProb(Q))$; then

\begin{align*}
& b'd'+c'(1-d')-a' \\
& =
(b + \sigma N(b(1-b)))(d + \sigma N(d(1-d)))
\\ & \quad +
(c + \sigma N(c(1-c)))(1 - d - \sigma N(d(1-d)))
\\ & \quad - (a+\sigma N(a(1-a))) \\
& = (bd+c(1-d)-a)
\\ & \quad + \sigma \left[
d N(b(1-b))\right. \\
&\quad\quad + (b-c) N(d(1-d)) \\
&\quad\quad + (1-d)N(c(1-c)) \\
&\left.\quad\quad - N(a(1-a))
\right]
\\ & \quad + O(\sigma^2) \\
\end{align*}

gives us a normal distribution with standard deviation

\begin{equation*}
\sigma\sqrt{
a(1-a)
+ d^2 b (1-b)
+ (1-d)^2 c (1-c)
+ (b-c)^2 d (1-d).
}
\end{equation*}

The violation metric is then:

\begin{equation*}
\frac{
|bd+c(1-d)-a|
}{
\sigma\sqrt{
a(1-a)
+ d^2 b (1-b)
+ (1-d)^2 c (1-c)
+ (b-c)^2 d (1-d).
}
}
\end{equation*}

\paragraph{Hyperparameters for hypothesis testing}
\label{cfc/app:frequentist-hyperparams}
Our goal is for the rejection criteria to be similar to the arbitrage violation metric in \Cref{cfc/app:arbitrage} on simple examples.
We choose $\gamma = 2.58$ for all checks, to ensure 99\%-confidence intervals for two-sided tests;
future work may consider using a different $\gamma$ for checks that require one-sided tests.
We pick $\sigma = 0.05$ (corresponding to $n = 400$ in \Cref{def:frequentist-consistency}).
The allowed violation threshold for all checks is then $\gamma \sigma = 0.129$.
For reference, a \NegChecker{} pair $(\Forecaster(P), \Forecaster(\neg P)) = (0.5, 0.59)$ has a violation metric of $0.128$, and would thus not be rejected as inconsistent.
This \rebuttal{closely} corresponds to the tolerance threshold of $10^{-2}$ of profit for the arbitrage metric, described in \Cref{cfc/sec:consistency-checks-intro}.

We pick $\betamin = 10^{-3}$ because LLM forecasters from \citet{halawiApproachingHumanLevelForecasting2024} answer with at most 3 digits of precision for events close to $0$ and $1$ in probability.



%\subsubsection{Tolerance scaling, [0, 1] space}
%
%\todo{this and the Adam section below need to be cleaned up}
%
%Assume that each event $P$ has a true probability value $\TrueProb(P)$. Take the null hypothesis:
%
%$\Forecaster$ produces a Gaussian estimate of each event $P$ centered at $\TrueProb(P)$, with variance $\sigma^2$ uniform across all events:
%
%\begin{align*}
%    \Forecaster(P) - \TrueProb(P) \sim \Normal(0, \sigma^2) \quad \text {independently for all events } P
%\end{align*}
%
%We take the violation metric and the corresponding threshold as to produce a hypothesis test against this:
%\begin{align*}
%   \Forecaster(P) + \Forecaster(\neg P) - 1
%   = \TrueProb(P) + \eps_1 + \TrueProb(\neg P) + \eps_2 - 1
%   = \eps_1 + \eps_2 \sim \Normal(0, 2\sigma^2)
%\end{align*}
%
%
%\subsubsection{Tolerance scaling, log space}
%For high or low probabilities, it is more principled to assume $\log \Forecaster$ produces a Gaussian estimate of the \emph{log prob}:
%
%\begin{align*}
%   \frac{\Forecaster(P)}{\TrueProb(P)} \sim \exp(\Normal(0, \sigma^2)) \quad
%    \text {independently for all events } P
%\end{align*}
%
%\dpaleka{Maybe it will be beneficial to have two sets of metrics reported, one behaving well for log error and one behaving well for linear error. Then make sure that those are both benevolent outside the intervals where they make sense, and reject consistency if at least one of them fails.}
%
%\todo{Adam's comment:  in practice log makes sense when we have a distribution bounded by 0 on one side and inf on the other..  Linear makes sense if left is unbounded and right is unbounded.  For prediction markets [0,1 bounds] something like logistic makes most sense}
%
%
%
%$$\exp(\varepsilon_1)\TrueProb(P)+\exp(\varepsilon_2)(1-\TrueProb(P))
%=\TrueProb(P)(\exp(\varepsilon_1)-\exp(\varepsilon_2))+\exp(\varepsilon_2) ???
%$$
%no wait the test statistic would actually be something else right not exactly $\Forecaster(P)+\Forecaster(\lnot P)-1$
%
%
%
%e.g., in the negation case log(p/(1-p))+log(q/(1-q)) has a distribution of
%
%$$\log\frac{X}{1-e^\varepsilon+X}$$
%
%where
%
%$X=e^{2\varepsilon}p^*(1-p^*)$

%\subsubsection{Asking LLM for confidence intervals and using Z-Scores as Violation Checker (Adam's methods)}\todo{this should be commented out no?}
%
%To determine the consistency of large language model (LLM) predictions, it is necessary to elicit not only prediction probabilities, but also confidence intervals. A more confident prediction should be associated with tighter violation tolerances, whereas looser confidence intervals indicate the model's understanding of its own uncertainty, allowing for more forgiving tolerances.
%
%As mentioned in \ref{cfc/sec:tolerance-scaling freq}, one method to ascertain a model's true prediction and confidence intervals is through a frequentist approach, which involves querying the LLM $n$ times and aggregating its outputs to estimate an average $\mu$ and a standard deviation $\sigma$ of its predictions. However, a key weakness of this approach is its assumption that the production of estimates from the LLM constitutes independent experiments. This assumption is flawed because, although LLMs are generally programmed to generate varied responses, the underlying data and architecture on which they are trained remain constant for each query. An analogous situation would be for a social scientist conducting a survey among 100 individuals regarding a prediction market event; there is a significant difference between querying 100 distinct individuals and querying a single individual 100 times.
%
%Consequently, the frequentist approach tends to underestimate the true $\sigma$ and confidence interval of the model. Since prediction events are singular and have not yet occurred, it is impossible to statistically construct accurate confidence intervals using past data. A feasible workaround is to directly ask the LLM model "What is your 1 $\sigma$ confidence interval for your previous prediction?" Similarly to human market participants and traders, this direct approach examines not only the accuracy and consistency of the model's predictions, but also its ability to comprehend uncertainty and probability ranges within a singular, unified framework.
%
%
%\subsubsubsection{Z-Score Tolerance Scoring}
%
%Another way to score consistency and tolerance is by utilizing the Z-score method.  Here we compare the amount of violation with the total expected confidence bounds that the model gave for its original predictions: $Z = \abs{\frac{\text{Violation}}{\sigma_{tot}}}$.  We can define $Z>1$ as a violation.
%
%Using the example of a modified negation checker, we separate $2\sigma^2$ to $\sigma_{orig}^2 \text{and} \sigma_{neg}^2 $, which are the associated stated confidence intervals that the LLM answers when asked for the original prediction and the negated prediction.  From there we calculate $Z = \frac{\eps_1 + \eps_2}{\sqrt{\sigma_{orig}^2 + \sigma_{neg}^2}}$.
%
%\begin{align*}
%   \Forecaster(P) + \Forecaster(\neg P) - 1
%   = \TrueProb(P) + \eps_1 + \TrueProb(\neg P) + \eps_2 - 1
%   = \eps_1 + \eps_2 \sim \Normal(0, {\sigma_{orig}}^2 + {\sigma_{neg}}^2)
%\end{align*}
%



%NegChecker: $\CCheck(\Forecaster(P),\Forecaster(\lnot P)):=\Forecaster(P)+\Forecaster(\lnot P)=1$; violation metric: $\Violation(\Forecaster(P),\Forecaster(\lnot P)) = |\Forecaster(P)+\Forecaster(\lnot P)-1|$.
%
%AndOrChecker: $\CCheck(\Forecaster(P),\Forecaster(Q),\Forecaster(P\land Q),\Forecaster(P\lor Q)):=\Forecaster(P)+\Forecaster(Q)=\Forecaster(P\lor Q)+\Forecaster(P\land Q)$; violation metric: $\Violation(\Forecaster(P),\Forecaster(Q),\Forecaster(P\land Q),\Forecaster(P\lor Q)) = |\Forecaster(P)+\Forecaster(Q)-\Forecaster(P\land Q)-\Forecaster(P\lor Q)|$.
%
%AndChecker. $\CCheck(\Forecaster(P),\Forecaster(Q),\Forecaster(P\land Q)):=\max(\Forecaster(P)+\Forecaster(Q)-1, 0) \le \Forecaster(P\land Q) \le \min(\Forecaster(P),\Forecaster(Q))$; violation metric: $\Violation(\Forecaster(P),\Forecaster(Q),\Forecaster(P\land Q)):=\max(\max(\Forecaster(P)+\Forecaster(Q)-1, 0)-\Forecaster(P\land Q), \Forecaster(P\land Q)-\min(\Forecaster(P),\Forecaster(Q))
%
%OrChecker: $\CCheck(\Forecaster(P),\Forecaster(Q),\Forecaster(P\lor Q)):=\max(\Forecaster(P),\Forecaster(Q))\le\Forecaster(P\lor Q)\le\min(1,\Forecaster(P)+\Forecaster(Q))$
%
%ButChecker: $\CCheck(\Forecaster(P),\Forecaster(\lnot P\land Q),\Forecaster(P\lor Q)):=\Forecaster(P\lor Q)=\Forecaster(P)+\Forecaster(\lnotP\land Q)$; violation metric: $\Violation(\Forecaster(P),\Forecaster(\lnot P\land Q),\Forecaster(P\lor Q)):=|\Forecaster(P)+\Forecaster(\lnot P\land Q)-\Forecaster(P\lor Q)|$.
%
%SymmetryAndChecker: $\CCheck(\Forecaster(P\land Q),\Forecaster(Q\land P)):=\Forecaster(P\land Q)=\Forecaster(Q\land P)$; violation metric: $\Violation(P\land Q,Q\land P)=|\Forecaster(P\land Q)-\Forecaster(Q\land P)|$
%SymmetryOrChecker: $\CCheck(\Forecaster(P\lor Q),\Forecaster(Q\lor P)):=\Forecaster(P\lor Q)=\Forecaster(Q\lor P)$; violation metric: $\Violation(P\lor Q,Q\lor P)=|\Forecaster(P\lor Q)-\Forecaster(Q\lor P)|$
%CondChecker: $\CCheck(\Forecaster(P),\Forecaster(Q|P),\Forecaster(P\land Q)); violation metric: $\Violation(\Forecaster(P),\Forecaster(Q|P),\Forecaster(P\land Q)):=\Forecaster(P)\Forecaster(Q|P)-\Forecaster(P\land Q)$.
%
%CondCondChecker: $\CCheck(\Forecaster(P), \Forecaster(Q|P), \Forecaster(R|P\land Q), \Forecaster(P\land Q\land R)):=\Forecaster(P)\Forecaster(Q|P)\Forecaster(R|P\land Q)=\Forecaster(P\land Q\land R)$; violation metric: \Violation(\Forecaster(P), \Forecaster(Q|P), \Forecaster(R|P\land Q), \Forecaster(P\land Q\land R))=\Forecaster(P)\Forecaster(Q|P)\Forecaster(R|P\land Q)-\Forecaster(P\land Q\land R)$
%
%ConsequenceChecker: (where $\Relation(P,Q):=P\implies Q$) $\CCheck(P,Q):=\Forecaster(P)\le\Forecaster(Q)$; violation metric: $\Violation(P,Q):=\max(\Forecaster(P)-\Forecaster(Q),0)$
%
%ParaphraseChecker (where $\Relation(P,Q):=P\iff Q$) $\CCheck(P,Q):=\Forecaster(P)=\Forecaster(Q)$; violation metric: $\Violation(P,Q)=|\Violation(P)-\Violation(Q)|$